\scanpage{183}
Soit $\mathcal M$ une catégorie de motifs géométriques, \inserted{de type \unsure{fini|lié|lisse},} sur le préschéma $S$ de type fini sur $\mathbf Z$. Pour simplifier, supposons que les degrés de $\mathcal M$ soient pairs (sinon, il sera possible, pour \deleted{la plupart des} \inserted{certaines} questions de distribution des Frobenius, de remplacer $\mathcal M$ par $\mathcal M^{\mathrm{pair}}$). Tout choix d'un pt fermé $s\in S$ permet donc, de définir des (une \unsure{collection|notion} de) ``$\mathbf R$-points Motiviques'' \deleted{\unsure{point de vue|de $\mathcal M$}''} \inserted{de $\mathcal M$}.

Pour fixer les idées, choisissons un pt $\mathbf R$-motivique de $\mathcal M$, d'où un gpe \inserted{physique} $G_{\mathbf R}$ sur $\mathbf R$.

Il y a lieu de \deleted{\unsure{dupliquer|d'expliciter}} \inserted{démultiplier} les \unsure{point|points} \inserted{fermés $s$} de $S$, en formant les couples $s'=(s,\unsure{i'|i|i_s|\xi})$, où $s\in S^0$, et
\margin{l'exposant $0$ désigne \unsure{des|les} pts \uline{fermés}}
$i_{\deleted{\unsure{s|x}}}\colon G_{m\,\unsure{\mathbf C|C}}-G(s)_{\mathbf C}$ est de type \deleted{\unsure{désigné|divisé}} préférentiel --- $G(s)$ désigne le gpe de Galois motivique associé à $\mathcal M_s$.

\deleted{Alors} Soit $S'$ l'ens de ces couples. A tout $s'\in S'$ est associé un $\mathbf R$-point motivique de $\mathcal M$, \deleted{d'où un} déterminé à isom près, d'où un $\unsure{\zeta|\xi}(s')\in H^1(\mathbf R,\deleted{\unsure{F|G}}G_{\mathbf R}^{\mathrm{or}})$, où $G^{\mathrm{or}}$ est le noyau de $G_{\mathbf R}^0\xrightarrow{\varepsilon}G_m$. De plus, si $G_{\mathbf R}^{\zeta(s')}$ est ``le'' gpe déduit de $G_{\mathbf R}$ par torsion par $\zeta(s')$, on obtient un hom
\[
u(s')\colon G(s')_{\deleted{\unsure{R|k}}\mathbf R}=G(s)_{\mathbf R}\longrightarrow G_{\mathbf R}^{\zeta(s')}
\]
compatible avec $\varepsilon_{\mathbf R},J_{\mathbf R}$, et déterminé en fait mod \deleted{\unsure{automorphisme|automorphismes}} \inserted{conjugaison} \deleted{\unsure{intérieur|interne|intérieurs}} par un élément de $G_{\mathbf R}^{\mathrm{or}}(\mathbf R)$.
\margin{182}
